\section{Introduction}
\label{sec:intro}

How many fish live in a lake? You cannot drain it to count them. Ecologists catch a net-full, tag each fish, and release them. A week later they catch another net-full and count how many carry tags. If a tenth of the second catch is tagged, the first catch was about a tenth of the lake. This is the capture--recapture method. Variants of it estimate wildlife populations, the number of species in a soil sample, and the number of people a census missed \citep{chao1987estimating,fienberg1972multiple,colwell2012models}.

Interpretability researchers face the same problem inside neural networks. A sparse autoencoder (SAE) rewrites a model's internal activations as a sparse combination of directions, called \emph{features}, many of which correspond to human-recognisable concepts \citep{cunningham2023sparse,bricken2023monosemanticity,templeton2024scaling}. A basic question has no accepted answer: \emph{how many features does a model have?} Dictionary size is chosen by hand. Larger dictionaries keep finding new latents \citep{gao2024scaling}. SAEs trained twice on the same data with different random seeds find partly different features \citep{paulo2025sparse}. Seed disagreement is usually treated as a nuisance.

We treat it as data. Each SAE seed is one net cast into the same lake. A feature found by several seeds has been recaptured, and a feature found by one seed was caught once. Whether some features are never found becomes a statistical question, the one capture--recapture answers. The analogy has one twist with no counterpart in fish. An SAE sometimes learns a direction that no other seed reproduces and that corresponds to no real feature. Ecologists know this failure from DNA-based surveys, where sequencing errors create ``species'' seen exactly once \citep{chiu2016estimating}.

\paragraph{Contributions.}
\begin{enumerate}\setlength\itemsep{1pt}
\item \textbf{A census method for features.} We cast SAE seed (dis)agreement as closed-population capture--recapture. We compare \NEstimators{} counts and estimators from ecology and statistics against planted ground truth, across \SynFamilies{} synthetic settings and \SynRuns{} trained SAEs (\S\ref{sec:method}).
\item \textbf{Naive counts fail, in a predictable direction.} Counting distinct features across seeds overstates the planted count by a median factor of \UnionRatioEq{} when the dictionary has as many latents as there are features, and by up to \UnionRatioMax$\times$ when it has more. Chao's classical estimator overstates it by up to $\ChaoRatioMax$. The cause is seed-specific junk latents: \JunkSingleOver{} of features caught once are junk when the dictionary is larger than the feature count (\S\ref{sec:results}).
\item \textbf{A singleton-free estimator that recovers the planted count.} It fits a bounded catch-probability mixture to features found by at least two seeds. Whenever the dictionary is at least as large as the true count, it is within \SfMaxErrComplete{} of the truth in all \NComplete{} settings. When the dictionary is too small, it tracks the features a dictionary of that size can reach, not the truth, so its rise and \emph{plateau} across dictionary sizes reveal the feature count. When planted features split into finer ones, the plateau appears only once the dictionary can hold the finest level (\S\ref{sec:hier}). It also predicts what held-out seeds will find.
\item \textbf{A census of Pythia-70M.} In the residual stream of Pythia-70M \citep{biderman2023pythia}, the estimated number of reproducible features grows from \RealSfOne{} to \RealSfEight{} as the dictionary grows from 1k to 8k latents, and growth slows without stopping. Fitted on four seeds, the census predicts what eight seeds reproduce to within \HeldRealSfMax. Five times more training data raises the count by a factor of \RealLongRatio{}, so all of these counts are lower bounds (\S\ref{sec:real}).
\end{enumerate}
